\scanpage{108}
\mspage{23}
% archive sheet number: 107
% source lines 1--3
\uline{Démonstration}. Soit $U_i=X-D_i$, $f_i:U_i\unsure{{-}|\longrightarrow}X$ l'inclusion, \deleted{de sorte qu'il \unsure{donne un}} \inserted{de sorte que} $f:U\unsure{{-}|\longrightarrow}X$ se factorise en

\diagram{\begin{tikzcd}[ampersand replacement=\&]
U \arrow[r,"g_i"] \arrow[rr,bend left=35,"f"] \& U_i \arrow[r,"f_i"] \& X.
\end{tikzcd}}

% source line 5
On en conclut un hom canonique

\[
\deleted{\unsure{(a)}}\qquad R^1 f_{i*}(\Lambda_{U_i})\unsure{{-}|\longrightarrow}R^1 f_*(\Lambda_U)
\]

% source line 7 and its correction
D'autre part, \deleted{(4.2.2) \unsure{nous dit et th. canonique}} \inserted{on a un isom (4.2.2)}

\[
R^1 f_{i*}(\Lambda_{U_i})\simeq\underline H^2_{D_i}(\Lambda_X)\simeq\deleted{\unsure{(\mu_n^{\otimes-1})_{D_i}}}\Lambda_{D_i}(-1),
\]

d'où un homomorphisme

\[
\unsure{\varphi_i|\varphi}:\deleted{\unsure{(\mu_n^{\otimes-1})_{D_i}}}\Lambda_{D_i}(-1)\longrightarrow R^1 f_*(\Lambda_U),
\]

et par suite un hom

\[
\deleted{\unsure{\text{avec}}}\qquad\varphi:\coprod_{1\leq i\leq r}\deleted{\unsure{(\mu_n^{\otimes-1})_{D_i}}}\Lambda_{D_i}(-1)\longrightarrow R^1 f_*(\Lambda_U).
\]

% source lines 14--15
Je dis que cet homomorphisme est un isomorphisme,

% The following block through the final words is crossed out by long diagonal strokes.
\deleted{Pour ceci, notons que $U=\bigcup_{1\leq i\leq r}U_i$, alors on a la suite spectrale de Leray}

\[
\deleted{R^{\unsure{*|p}}f_*(\Lambda_U)\Longleftarrow E^{pq}_2=H^p((i_0\cdots i_p)\longrightarrow R^q f_{i_0\cdots i_p*}(\Lambda_{U_{i_0\cdots i_p}})\unsure{)})}
\]

\deleted{où $U_{i_0\cdots i_p}=U_{i_0}\cap\cdots\cap U_{i_p}$, et où \deleted{\unsure{$f$.}} $f_{i_0\cdots i_p}:U_{i_0\cdots i_p}\longrightarrow X$ est l'inclusion canonique.}

\deleted{On a évidemment $R^0 f_{i_0\cdots i_p*}(\Lambda_{U_{i_0\cdots i_p}})\simeq\Lambda_X$, donc $E^{p0}_2=0$ pour tout $p\ne 0$, d'où résulte que la suite spectrale donne un isomorphisme}

\[
\deleted{\deleted{\unsure{E_2}}\quad R^1 f_*(\Lambda_U)\xrightarrow{\sim}E^{01}_2\subset\coprod_{1\leq i\leq r}R^1 f_{i*}(\Lambda_{U_i})}
\tag*{\deleted{(***)}}
\]

\deleted{Or on constate tout de suite que le}

% small drawing in the lower-left margin
\figure{\begin{tikzpicture}[x=0.6cm,y=0.6cm]
\draw (-0.65,0)--(3.7,0);
\draw (0,-0.6)--(0,3.4);
\node at (0,0) {$\times$};
\node at (0,0.48) {$\times$};
\node at (0.9,0.52) {$\times$};
\node at (1.7,0.57) {$\times$};
\node at (2.55,0.68) {$\times$};
\end{tikzpicture}}
